\documentclass[11pt]{article}

\usepackage[a4paper,margin=1in]{geometry}
\usepackage{amsmath,amssymb,amsthm,mathtools}
\usepackage{enumitem}
\usepackage{hyperref}

\title{A Geometric Interpretation of the BDV Selection Principle\newline\large Toward a Flow Picture for Quantitative Faber--Krahn}
\author{}
\date{}

\newcommand{\Om}{\Omega}
\newcommand{\R}{\mathbb{R}}
\newcommand{\Sph}{\mathbb{S}}
\newcommand{\A}{\mathcal A}
\newcommand{\D}{\mathcal D}
\newcommand{\fint}{\mathop{\mathchoice{\vcenter{\hbox{\Large$\int$}}\kern-1.4em\vcenter{\hbox{\scriptsize$-$}}}{\vcenter{\hbox{\large$\int$}}\kern-1.1em\vcenter{\hbox{\scriptsize$-$}}}{\vcenter{\hbox{$\int$}}\kern-.85em\vcenter{\hbox{\scriptsize$-$}}}{\vcenter{\hbox{$\int$}}\kern-.7em\vcenter{\hbox{\tiny$-$}}}}\nolimits}

\begin{document}

\maketitle

\section{Guiding philosophy}

For the quantitative isoperimetric inequality, one can think of the hidden mechanism as a geometric motion toward the extremal set, namely the ball. Depending on the proof, this motion may be implemented by symmetrization, transport, or another comparison deformation. The important conceptual point is that there is a process
\[
    E \rightsquigarrow B
\]
along which the relevant quotient improves. Stability then comes from quantifying the amount by which the quotient improves: if almost no improvement occurs, the original set must already be close to a ball.

For Faber--Krahn, the analogous mechanism should not be a raw transport of volume. The relevant object is not only the domain $\Om$, but also the first eigenfunction $\varphi_\Om$. The eigenfunction carries the geometry of the domain through its level sets and through the boundary pressure $(\partial_\nu \varphi_\Om)^2$.

Thus the proposed geometric picture is:
\[
    \text{move the boundary so as to equalize eigenfunction pressure.}
\]
The ball is characterized by the fact that the first eigenfunction has constant normal derivative on the boundary. Therefore the ball is the equilibrium of a pressure-equalizing motion.

\section{The explicit local flow}

Let $\Om_t \subset \R^n$ be a smooth family of domains with fixed volume, and let $\varphi_t$ be the positive first Dirichlet eigenfunction, normalized by
\[
    \int_{\Om_t} \varphi_t^2 = 1.
\]
The Hadamard shape derivative gives, for a normal deformation with velocity $V_n$,
\[
    \frac{d}{dt}\lambda_1(\Om_t)
    = - \int_{\partial \Om_t} (\partial_\nu \varphi_t)^2 V_n \, d\sigma.
\]
The volume constraint is
\[
    \frac{d}{dt}|\Om_t| = \int_{\partial \Om_t} V_n \, d\sigma = 0.
\]
Hence the natural volume-preserving steepest descent is
\[
    V_n
    = (\partial_\nu \varphi_t)^2
    - \fint_{\partial \Om_t} (\partial_\nu \varphi_t)^2 \, d\sigma .
\]
Along this flow,
\[
\begin{aligned}
    \frac{d}{dt}\lambda_1(\Om_t)
    &= -\int_{\partial \Om_t}
    \left((\partial_\nu \varphi_t)^2
    - \fint_{\partial \Om_t}(\partial_\nu \varphi_t)^2\,d\sigma\right)^2 d\sigma \\
    &\leq 0.
\end{aligned}
\]
So the monotone quantity is the first eigenvalue at fixed volume, and the dissipation is
\[
    \D(\Om)
    :=
    \int_{\partial \Om}
    \left((\partial_\nu \varphi_\Om)^2
    - \fint_{\partial \Om}(\partial_\nu \varphi_\Om)^2\,d\sigma\right)^2 d\sigma .
\]
This is the Faber--Krahn analogue of the perimeter decrease in a geometric proof of the isoperimetric inequality.

The equilibrium condition is
\[
    (\partial_\nu \varphi)^2 = \text{constant on } \partial \Om.
\]
By Serrin's rigidity theorem, a smooth bounded connected domain satisfying the overdetermined problem is a ball. Thus the flow is geometrically trying to drive the domain toward the unique pressure-balanced shape.

\section{Geometric meaning}

The boundary pressure $(\partial_\nu\varphi)^2$ measures how much the eigenvalue benefits from expanding the boundary at a given point. If this pressure is larger than average, the flow pushes the boundary outward. If it is smaller than average, the flow pulls the boundary inward. Therefore the motion redistributes the boundary until all boundary points have equal spectral marginal value.

This is analogous to curvature equalization in the isoperimetric problem. For isoperimetry, the ball is characterized by constant mean curvature. For Faber--Krahn, the ball is characterized by constant eigenfunction pressure. Hence:
\[
    \text{mean-curvature equalization}
    \quad \leftrightarrow \quad
    \text{eigenfunction-pressure equalization}.
\]

Rough boundary features are expected to disappear under this motion. Thin tentacles, spikes, fjords, and high-frequency oscillations are spectrally inefficient: they contribute volume and asymmetry but contribute little useful eigenfunction mass. The eigenfunction acts as a harmonic low-pass filter. The pressure flow erases boundary features that are not supported by the eigenfunction.

\section{Relation with rearrangement and level sets}

The classical Schwarz symmetrization proof of Faber--Krahn can be viewed as a level-set version of this philosophy. One does not directly transport $\Om$ to a ball; instead one rearranges the level sets of the eigenfunction:
\[
    \{\varphi_\Om > t\} \rightsquigarrow \{\varphi_\Om^* > t\}.
\]
The P\'olya--Szeg\H{o} inequality says that the Dirichlet energy decreases under this rearrangement while the $L^2$ norm is preserved. Hence the Rayleigh quotient decreases.

The quantitative idea would be: if the Rayleigh quotient barely decreases, then for many important levels $t$, the sets $\{\varphi_\Om>t\}$ nearly saturate the isoperimetric inequality. This suggests that most relevant level sets are close to balls. The difficulty is to recover good control of the boundary level $\{\varphi_\Om>0\}=\Om$, especially for rough sets. This is where the BDV selection principle enters.

\section{BDV as an implicit version of the flow}

The Brasco--De Philippis--Velichkov proof does not explicitly run the pressure-equalizing flow. Instead, the selection principle can be understood as an implicit minimizing-movement step for such a flow.

Starting from a rough almost minimizer, BDV introduce a penalized problem that selects a better representative among nearby competitors with comparable asymmetry scale. The selected set satisfies an Euler--Lagrange/free-boundary condition of Bernoulli type:
\[
    |\nabla u|^2 = \Lambda + \text{small forcing on }\partial U.
\]
This says that the selected set is an almost-equilibrium for the pressure law. Thus the selection principle is geometrically doing the following:
\[
    \text{rough almost minimizer}
    \quad \rightsquigarrow \quad
    \text{smooth almost pressure-balanced set}.
\]
It replaces the missing explicit flow by its endpoint. The rough set is not directly analyzed; instead it is projected onto a nearby free-boundary equilibrium where regularity theory and local spectral analysis become available.

In this sense, BDV's selection principle is not merely a technical compactness device. It is a geometric relaxation mechanism: it discards spectrally inefficient boundary features while keeping the relevant asymmetry scale.

\section{The local spectral picture near the ball}

Once the selected set is smooth and close to a ball, one can write its boundary as a normal graph
\[
    \partial \Om = \{(1+\rho(\theta))\theta : \theta \in \Sph^{n-1}\}.
\]
The volume constraint removes the spherical harmonic mode $\ell=0$, and centering the barycenter removes the translation modes $\ell=1$. The first nontrivial modes are therefore $\ell\geq 2$.

Near the ball, one expects expansions of the form
\[
    \lambda_1(\Om)-\lambda_1(B)
    \simeq Q(\rho,\rho),
\]
where $Q$ is positive definite on the modes $\ell\geq 2$. Meanwhile, the pressure defect
\[
    (\partial_\nu\varphi_\Om)^2
    - \fint_{\partial\Om}(\partial_\nu\varphi_\Om)^2\,d\sigma
\]
linearizes to an elliptic operator on $\rho$ whose kernel consists precisely of the forbidden volume and translation modes. Thus, after imposing volume and barycenter constraints, one should have a spectral gap estimate
\[
    \|\rho\|_{H^{1/2}(\Sph^{n-1})}^2
    \lesssim
    \left\|
    (\partial_\nu\varphi_\Om)^2
    - \fint_{\partial\Om}(\partial_\nu\varphi_\Om)^2\,d\sigma
    \right\|_{L^2(\partial\Om)}^2.
\]
Combining this with the expansion of the eigenvalue gives the local stability mechanism.

Finally,
\[
    \A(\Om)^2
    \lesssim \|\rho\|_{L^2(\Sph^{n-1})}^2
    \lesssim \|\rho\|_{H^{1/2}(\Sph^{n-1})}^2,
\]
so control of the pressure defect implies control of the Fraenkel asymmetry.

\section{Summary}

The proposed geometric interpretation is:
\[
\boxed{
\text{Quantitative Faber--Krahn is governed by eigenfunction-pressure equalization.}
}
\]
The explicit formal flow is
\[
\boxed{
    V_n
    = (\partial_\nu\varphi)^2
    - \fint_{\partial\Om}(\partial_\nu\varphi)^2\,d\sigma .
}
\]
Along this flow, at least formally for smooth domains,
\[
\boxed{
    \frac{d}{dt}\lambda_1(\Om_t)
    = -\D(\Om_t) \leq 0.
}
\]
The ball is the unique equilibrium by Serrin's theorem. The local stability mechanism is the spectral gap of the linearized pressure map around the ball, after removing the volume and translation modes.

BDV's selection principle can then be read as an implicit rigorous substitute for the flow in the rough global regime:
\[
\boxed{
    \text{selection principle}
    \approx
    \text{implicit pressure-flow relaxation to a smooth almost-equilibrium.}
}
\]
The proof does not explicitly transport the original domain to the ball. Instead, it replaces the original domain by a nearby smooth pressure-balanced representative, and then uses the rigidity of the ball as the unique equilibrium together with local spectral coercivity.

\end{document}
